\documentclass[11pt]{letter}
\usepackage[margin=1in]{geometry}
\usepackage{hyperref}

\signature{Elliot Tower \\ Independent researcher \\ \texttt{elliot@elliottower.ai}}

\begin{document}

\begin{letter}{Editorial Board \\ \emph{GigaScience} \\ Oxford University Press}

\opening{Dear Editors,}

I am submitting the manuscript ``Edge-Level Discrete Curvature Decomposes Bridge Structure Across Disease Networks'' for consideration as a Research Article in \emph{GigaScience}.

Ollivier-Ricci curvature has been applied to brain, cancer, and autism networks, but every biomedical application to date reports curvature at the node level without a null model separating curvature from degree. We show that node-level curvature is a degree artifact ($r = -0.71$) and present an edge-level analysis with degree-preserving and weight-permutation nulls that recovers genuine bridge structure. The method replicates across seven independent graphs spanning five data types and three disease domains: in every graph, the lowest-curvature node is a recognized diagnostic bridge entity. Applied to a psychiatric mechanism network, disconfirmed claims occupy significantly lower-curvature positions than higher-verdict claims ($p = 0.007$, $d = 0.89$), and among five edge-level methods, only curvature achieves this discrimination.

\emph{GigaScience} is a good fit for three reasons:

\begin{enumerate}
\item \textbf{Data-intensive, cross-domain scope.} The analysis spans seven graphs constructed from genetic correlations, shared GWAS loci, epidemiological odds ratios, hazard ratios, and literature curation across psychiatric, autoimmune, and cardiometabolic disease. This breadth suits a venue that values large-scale, multi-data-type analyses.

\item \textbf{Methodological contribution with immediate reuse.} The degree-preserving and weight-permutation null models are portable to any biomedical network where curvature is computed. The paper provides a concrete correction for a pervasive confound in an active area of network biology.

\item \textbf{Open data and open code.} All analysis scripts, the mechanism catalog, precomputed curvatures, permutation distributions, and MR results are available under an MIT license at \url{https://github.com/elliottower/psychiatric-comorbidity-curvature}.
\end{enumerate}

I used AI assistants (Claude Code by Anthropic; Perplexity AI) for text editing, code drafting/debugging, and literature review assistance in constructing the mechanism catalog. I reviewed and verified all outputs and take full responsibility for the content. Full details are disclosed in the manuscript.

I confirm that this manuscript has not been published previously and is not under consideration at another journal. I have no competing interests and this work received no external funding.

\closing{Sincerely,}

\end{letter}
\end{document}
